% ---- pg-427 （书页 415） ----

\[
\begin{array}{l}
\textbf{变\quad 量}\\[4pt]
\text{资产历史收益率,HRA}\\
\text{股本历史收益率,HRE}\\
\text{销售历史收益率,HRS}\\
\text{资产重置收益率,RRA}\\
\text{股本重置收益率,RRE}\\
\text{销售重置收益率,RRS}\\
\text{市场 Q 比率,Q}\\
\text{市场相对溢价,REV}
\end{array}
\qquad
\left[\begin{array}{rrrrrrrr}
\text{HRA}&\text{HRE}&\text{HRS}&\text{RRA}&\text{RRE}&\text{RRS}&\text{Q}&\text{REV}\\
1.000& & & & & & & \\
0.738&1.000& & & & & & \\
0.731&0.520&1.000& & & & & \\
0.828&0.688&0.652&1.000& & & & \\
0.681&0.831&0.513&0.887&1.000& & & \\
0.712&0.543&0.826&0.867&0.692&1.000& & \\
0.625&0.322&0.579&0.639&0.419&0.608&1.000& \\
0.604&0.303&0.617&0.563&0.352&0.610&0.937&1.000
\end{array}\right]
\]

得到对 $m=3$ 因子模型因子载荷的旋转主成分估计如下:

\begin{mvtable}
\small
\begin{tabular}{l|ccc}
\toprule
变\quad 量 & \multicolumn{3}{c}{因子载荷估计}\\
\cmidrule(l){2-4}
& $F_1$ & $F_2$ & $F_3$\\
\midrule
HRA & 0.433 & 0.612 & 0.499\\
HRE & 0.125 & 0.892 & 0.234\\
HRS & 0.296 & 0.238 & 0.887\\
RRA & 0.406 & 0.708 & 0.483\\
RRE & 0.198 & 0.895 & 0.283\\
RRS & 0.331 & 0.414 & 0.789\\
Q & 0.928 & 0.160 & 0.294\\
REV & 0.910 & 0.079 & 0.355\\
\midrule
解释全部方差的累积比例 & 0.287 & 0.628 & 0.908\\
\bottomrule
\end{tabular}
\end{mvtable}

(a) 利用因子载荷估计, 确定特殊方差和共性方差.

(b) 确定残差矩阵 $\mathbf{R}-\hat{\mathbf{L}}_z\hat{\mathbf{L}}_z'-\hat{\boldsymbol{\Psi}}_z$. 已知此信息和上面表格中的解释总方差的累积比例, $m=3$ 的因子模型对这些数据是否恰当?

(c) 假设小于 0.4 的载荷估计是小的, 解释这三个因子. 例如, 会出现这种情况吗: 市场价值测量值表明的利润与会计测量值表明的利润不同? 你能将会计历史利润的测量值从会计重置测量中分出来吗?

\noindent 9.16\quad 证明由\hyperref[eq:9-50]{式(9-50)}构造的因子得分有样本均值向量 $\mathbf{0}$ 和零样本协方差.

\noindent 9.17\quad 参考例 9.12, 利用该例中的信息评估 $(\hat{\mathbf{L}}_z'\hat{\boldsymbol{\Psi}}_z^{-1}\hat{\mathbf{L}}_z)^{-1}$. 注意: 令 $\hat{\boldsymbol{\Psi}}_z$ 的第四对角元素为 0.01, 这样可以确定 $\hat{\boldsymbol{\Psi}}_z^{-1}$. 对标准化股价观测构造因子得分的回归和广义最小二乘法将得出几乎相同的结果吗? 提示: 参见\hyperref[eq:9-57]{式(9-57)}和随后的讨论.

\mvsub{以下练习需使用计算机.}

\noindent 9.18\quad 参见练习 8.16 的捕鱼数据
(a) 仅用测量值 $x_1-x_4$, 求出 $m=1$ 和 $m=2$ 的因子模型的主成分解.
(b) 仅用测量值 $x_1-x_4$, 求出 $m=1$ 和 $m=2$ 的因子模型的极大似然解.
(c) 旋转(a)和(b)的解. 比较这些解, 并加以评论. 解释每一个因子.
(d) 用测量值 $x_1-x_6$ 执行因子分析, 求因子的合理个数 $m$, 并比较旋转之后的主成分和
